\section{Core fields and layered architecture}

\subsection{Metric, condensate, and finite-density frame}

The gravitational variable is the spacetime metric $g_{\mu\nu}$.  The
candidate plenum order parameter is complex,
\begin{equation}
 \Phi=\frac{\rho}{\sqrt{2}}e^{i\Theta},
 \label{eq:complex-order}
\end{equation}
so that amplitude fluctuations, phase excitations, zeros of the order
parameter and winding sectors can be represented within one microscopic
language.  UVIR-001 rejects the minimally kinetic analytic candidate as the
direct origin of the spatial cubic force operator, but it does not remove
these density, circulation, defect and topology roles.

A unit timelike vector $U^\mu$ defines the preferred frame,
\begin{equation}
 U^\mu U_\mu=-1,
 \qquad
 h^{\mu\nu}=g^{\mu\nu}+U^\mu U^\nu.
 \label{eq:projector}
\end{equation}
UVIR-003 Stage A provisionally chooses $U^\mu$ as an independent constrained
dynamical field.  The condensate current may be aligned with it through an
interaction energy, but is not algebraically identified with it.  This avoids
double-counting a derived condensate velocity and an independent aether field.
The choice remains provisional until the coupled constraint and stability gate
closes.

\subsection{Separate preferred-frame force scalar and matter}

The low-acceleration force mediator is denoted $\psi$.  Under the selected
UVIR-002 architecture it is a separate preferred-frame scalar, not the phase
of $\Phi$ relabelled at low energy.  Independent temporal and spatial
invariants can therefore be formed from $U^\mu\nabla_\mu\psi$ and
$h^{\mu\nu}\nabla_\mu\psi\nabla_\nu\psi$.  The cubic spatial operator is a
conditional infrared EFT input; it has not been matched to a microscopic
completion.

Matter fields are collected in $\Psi_m$.  The matter--force interaction must
be derived once, then used consistently in both the scalar source and the
matter stress-energy exchange.  No coefficient inferred from a convenient
normalization is promoted to an observable coupling before MAT-001.

\subsection{Topology, reservoir, and non-equilibrium sectors}

Spatial slices are assigned compact $\Tthree$ boundary conditions with lengths
$L_i$.  Their ratios are shape moduli.  The global topology controls allowed
modes and winding; it does not directly fix every local Wilson coefficient.

An explicit reservoir sector $R$ permits the observable matter--plenum system
to be open while the complete covariant system remains conserved.  A wake or
memory variable $W$ and a shape sector are permitted research sectors, not
assumed completed degrees of freedom.

The schematic inventory is
\begin{equation}
 S_{\rm complete}=S_{\rm EH}+S_\Phi+S_U+S_{\rm align}+S_\psi
 +S_m+S_R+S_{\rm int}+S_W,
 \label{eq:inventory}
\end{equation}
where $S_W$ is included only if a causal and stable wake description is found.
The displayed sectors define a research architecture, not a completed
microscopic action.  Effective operators must be matched and overlapping UV
and IR descriptions must not be counted as independent copies of the same
degree of freedom.

\subsection{Architecture versus prediction}

Equation~\eqref{eq:inventory} is a sector map.  In particular,
\begin{equation}
 \text{energy throughput}
 \not\Rightarrow
 \text{anisotropic pressure}
\end{equation}
without a shape-dependent conversion law, and
\begin{equation}
 \text{compact topology}
 \not\Rightarrow
 \cproj=\frac23 \quad\text{or}\quad \frac{H_t}{H_p}=\frac{13}{12}
\end{equation}
without separate matching and backreaction calculations.  Likewise, a
positive frozen-background dispersion is not a proof of coupled stability,
causality or weak coupling.
